% Lecture example set: SC3000-L05
% Source: CZ3005 Module 5 pp. 23--27
% Human verified: Yes

\clearpage
\exercisemeta
  {SC3000-L05-EX}
  {2026-09-14}
  {Module 5 pp.~23--27}
  {公式、数值与讲义内部参数差异已核对}
  {已确认（2026-09-14）}

\exercisequestion{SC3000-L05-E01}{五房间 Q-Learning 的奖励传播}

\textbf{对应知识点：}
\begin{itemize}
  \item \relatedtopic{SC3000-L05-T01}{Temporal Difference 与 Q-Learning}
  \item \relatedtopic{SC3000-L05-T02}{Off-Policy Control 与算法边界}
\end{itemize}

\subsubsection*{题目}

讲义把房间 $0,\ldots,5$ 视为 states，合法房门视为 actions；进入目标房间 $5$ 的即时 reward 为 $100$，其他合法移动 reward 为 $0$。到达 $5$ 后 episode 结束。初始化 $Q_0(s,a)=0$，并取 $\alpha=0.01,\gamma=0.99$。

\begin{figure}[htbp]
  \centering
  \begin{tikzpicture}[
  >=Latex,
  state/.style={circle,draw=noteBlue,fill=noteBlue!8,thick,minimum size=0.9cm},
  goal/.style={circle,draw=ntuRed,fill=ntuRed!10,thick,minimum size=0.9cm},
  ordinary/.style={<->,thick,draw=gray!80},
  reward/.style={->,very thick,draw=ntuRed},
  lab/.style={font=\scriptsize,fill=white,inner sep=1pt}
]
  \node[state] (s1) at (0,2.2) {$1$};
  \node[state] (s3) at (0,0.8) {$3$};
  \node[state] (s4) at (0,-0.6) {$4$};
  \node[state] (s2) at (-2.2,0.8) {$2$};
  \node[state] (s0) at (-2.2,-0.6) {$0$};
  \node[goal]  (s5) at (2.4,0.8) {$5$};

  \draw[ordinary] (s1) -- node[lab,left]{$r=0$} (s3);
  \draw[ordinary] (s3) -- node[lab,left]{$r=0$} (s4);
  \draw[ordinary] (s2) -- node[lab,above]{$r=0$} (s3);
  \draw[ordinary] (s0) -- node[lab,below]{$r=0$} (s4);
  \draw[reward] (s1) -- node[lab,above,sloped]{$r=100$} (s5);
  \draw[reward] (s4) -- node[lab,below,sloped]{$r=100$} (s5);
  \node[align=left,font=\small,right=0.55cm of s5] {terminal\\state};
\end{tikzpicture}

  \caption{讲义五房间环境的 episodic 版本；红色边进入 terminal state 5 并获得 reward 100。}
  \label{fig:sc3000-l05-five-room}
\end{figure}

计算第一次 transition $1\to5$ 后的 $Q(1,5)$；随后假设 $Q(1,5)=1$，计算一次 transition $3\to1$ 后的 $Q(3,1)$。说明奖励如何向前传播并提取 policy。

\begin{checkedsolution}
由于 $5$ 是 terminal state，第一次更新没有 future-value term：
\[
  Q_{\mathrm{new}}(1,5)
  =0+0.01(100-0)
  =\boxed{1}.
\]
随后从 $3$ 进入 $1$ 时即时 reward 为零，而 $\max_aQ(1,a)=Q(1,5)=1$，故
\[
  Q_{\mathrm{new}}(3,1)
  =0+0.01\left(0+0.99\times1-0\right)
  =\boxed{0.0099}.
\]
继续采样后，目标 reward 会沿合法路径逐步传播到离目标更远的 state-action pairs。训练后在每个非终止 state 选择
\[
  \pi(s)=\arg\max_aQ(s,a).
\]
实际收集经验仍可采用 $\varepsilon$-greedy behavior policy，以免尚未尝试的 actions 永远得不到更新。
\end{checkedsolution}

\subsubsection*{讲义参数核对}

Module 5 p.~25 写 $\gamma=0.99$，但 p.~27 的最终表出现 $100,80,64,51$；这些数值满足
\[
  100,\quad0.8(100),\quad0.8^2(100),\quad0.8^3(100)\approx51,
\]
所以该表实际对应 $\gamma=0.8$。若使用 $\gamma=0.99$，相邻距离的值应接近 $100,99,98.01,97.03$。此外，讲义 reward matrix 仍给 state $5$ 保留 outgoing actions；若严格按“到达 5 即结束”的 episodic 定义，这些 actions 不再执行。做题时应以题目明确给出的 terminal condition 与 $\gamma$ 为准。

\textbf{检查点：}\par
terminal transition 的 target 只有即时 reward；非终止 transition 才加入 $\gamma\max_{a'}Q(s',a')$。
